\documentclass[11pt]{article}

\usepackage[T1]{fontenc}
\usepackage{amsmath,amssymb,amsthm}
\usepackage{hyperref}

\title{Relative-fit and nested difference tests for estimated-weight
  moment-quadratic estimators}
\author{}
\date{\today}

\newcommand{\E}{\mathbb{E}}
\newcommand{\tr}{\operatorname{tr}}
\newcommand{\Var}{\operatorname{Var}}
\newcommand{\Cov}{\operatorname{Cov}}
\newcommand{\vech}{\operatorname{vech}}
\newcommand{\diag}{\operatorname{diag}}
\newcommand{\argmin}{\operatorname*{arg\,min}}
\newcommand{\argmax}{\operatorname*{arg\,max}}
\newcommand{\plim}{\operatorname*{plim}}
\newcommand{\T}{^{\mathsf{T}}}
\newcommand{\Wr}{W_{\!r}}
\newcommand{\rs}{r_{*}}
\newcommand{\ths}{\theta_{*}}

\theoremstyle{plain}
\newtheorem{proposition}{Proposition}
\newtheorem{theorem}{Theorem}
\newtheorem{lemma}{Lemma}
\newtheorem{corollary}{Corollary}
\theoremstyle{remark}
\newtheorem*{remark}{Remark}

\begin{document}
\maketitle

\section*{Purpose}

The estimated-weight sandwich (companion note
\texttt{weighted\_moment\_ij\_grid}) corrects the parameter covariance by
adding a weight-derivative term. This note asks whether the same correction
extends to model-comparison tests, and whether the nested difference test
follows as a corollary. Three answers, stated up front.

\begin{enumerate}
\item The correction is proportional to the pseudo-true residual $\rs$
  (Proposition~\ref{prop:delin}). It is identically zero for the
  goodness-of-fit null, where the focal model is assumed correct, and present
  for every test whose null does not assume correctness.
\item For distinguishable models the relative-fit (Vuong) variance is a clean
  corollary: $\omega^{2}=c\T\Gamma c$ with the weight channel carried by an
  explicit term in $c$ (Proposition~\ref{prop:vuong}).
\item The nested, degenerate case is \emph{not} a clean corollary. The weight
  channel and the bread curvature both deform Satorra's projection, so the
  naive substitution $\Gamma\mapsto(W-\Wr)\Gamma(W-\Wr)\T$ into the $U$-matrix
  is wrong (Proposition~\ref{prop:nested}). It needs its own derivation.
\end{enumerate}

\section{Recap of the corrected influence}
\label{sec:recap}

The estimator minimises $F_{N}(\theta)=\tfrac12 r(\theta)\T\widehat{W}r(\theta)$
with residual $r(\theta)=\eta(\theta)-u$, Jacobian $\Delta=\partial\eta/\partial
\theta\T$, and weight $\widehat{W}=W(u)$ built from the sample moments $u$. Write
$g_{i}$ for the moment influence function, $\Gamma=\E[g_{i}g_{i}\T]$, and $\ths$,
$\rs=\eta(\ths)-u_{0}$ for the pseudo-true value and residual. The estimator is
asymptotically linear with influence function $\phi_{i}=H^{-1}\Delta\T
v_{i}$, where
\begin{equation}
  H = A + C,\quad A=\Delta\T W\Delta,\quad
  C=\sum_{k}(W\rs)_{k}\nabla^{2}\eta_{k}(\ths),\quad
  v_{i}=Wg_{i}-\dot W_{i}\,\rs,
  \label{eq:recap}
\end{equation}
where $\dot W_{i}$ is the casewise influence of the weight. When the weight is a
function of the modelled moments $u$, $\dot W_{i}\rs=\Wr g_{i}$ with
$\Wr=[\,\partial W/\partial u_{1}\,\rs,\ \dots,\ \partial W/\partial u_{m}\,
\rs\,]$ the residual-contracted weight derivative, so $v_{i}=(W-\Wr)g_{i}$ and
the corrected meat is $\Omega=(W-\Wr)\Gamma(W-\Wr)\T$. In general $\Omega=
\Var(v_{i})$ is formed from the casewise rows; the matrix form is the special
case used in Sections~\ref{sec:gof}--\ref{sec:nested}.

\section{The residual governs the correction}
\label{sec:delin}

\begin{proposition}
\label{prop:delin}
Every estimated-weight correction in \eqref{eq:recap}, namely the curvature
term $C$ and the weight term $\Wr$, is linear in the pseudo-true residual
$\rs$. Hence both vanish exactly when $\rs=0$. For a test whose null fixes the
focal model as correctly specified the correction is identically zero; for a
test whose null leaves the focal model misspecified the correction is present.
\end{proposition}

\begin{proof}
$C$ contracts $W\rs$ and $\Wr$ contracts $\rs$ columnwise, so both are $O(\lVert
\rs\rVert)$ and vanish at $\rs=0$. The goodness-of-fit null sets $\eta(\ths)=
u_{0}$, hence $\rs=0$. The relative-fit and the not-assuming-correct nested
nulls leave $\rs\neq0$.
\end{proof}

A second, sharper reason closes the goodness-of-fit case. At that null the
fitted residual is $r(\widehat\theta)=O_{p}(N^{-1/2})$, so the weight
fluctuation enters the residual only through $(\widehat W-W)r(\widehat\theta)=
O_{p}(N^{-1})$, which is negligible. The estimated weight is therefore
asymptotically irrelevant to the fit statistic, which is why the fixed-weight
projection theory of \cite{satorra1989} and the scaled statistic implemented in
practice are already correct under the null.

\section{Goodness of fit: the projection, unchanged}
\label{sec:gof}

For a fixed weight, expand the fitted residual at the goodness-of-fit null. With
$H=A$ and $\Wr=0$ there,
\[
  \sqrt{N}\,r(\widehat\theta)=-M\sqrt{N}(u-u_{0})+o_{p}(1),
  \qquad M=I-\Delta A^{-1}\Delta\T W,
\]
so $\Var(\sqrt N\,r(\widehat\theta))=M\Gamma M\T$ and the fit statistic
$N\,r\T W r$ converges to $\sum_{j}\lambda_{j}\chi^{2}_{1,j}$ with $\lambda_{j}=
\operatorname{eig}(U\Gamma)$ and $U=WM=W-W\Delta(\Delta\T W\Delta)^{-1}\Delta\T
W$ \cite{satorra1989}. When $W=\Gamma^{-1}$ the product $U\Gamma$ is idempotent
and the law collapses to $\chi^{2}_{\mathrm{df}}$; for $W\neq\Gamma^{-1}$ (DWLS,
ULS) it is a genuine mixture, the reason the scaled statistic
\cite{satorrabentler2001} exists. By Proposition~\ref{prop:delin} the estimated
weight contributes nothing here.

\section{Relative fit: a clean corollary}
\label{sec:vuong}

Compare two structural models $A$ and $B$ on the same indicators. The weight
$\widehat W=W(u)$ is a function of the sample moments alone, so it is shared,
$W_{A}=W_{B}=W$. Let $D=F_{A}(\widehat\theta_{A})-F_{B}(\widehat\theta_{B})$ be
the sample discrepancy difference and $\rs^{A},\rs^{B}$ the two pseudo-true
residuals.

\begin{proposition}
\label{prop:vuong}
The discrepancy difference is asymptotically linear, $\sqrt{N}\,D=N^{-1/2}\sum_{i}
c\T g_{i}+o_{p}(1)$, with
\begin{equation}
  c = -W(\rs^{A}-\rs^{B}) + \tfrac12(b^{A}-b^{B}),
  \qquad
  b^{A}_{l}=(\rs^{A})\T\frac{\partial W}{\partial u_{l}}\,\rs^{A},
  \label{eq:cvec}
\end{equation}
and $b^{B}_{l}$ defined likewise. For distinguishable models $\omega^{2}=c\T
\Gamma c>0$ and $\sqrt{N}\,D/\widehat\omega\to\mathcal N(0,1)$.
\end{proposition}

\begin{proof}
View $F_{A}$ as a functional of the sampling distribution through $u$, the
minimiser $\widehat\theta_{A}(u)$, and the weight $W(u)$. The minimiser is
interior, so the envelope theorem removes the $\widehat\theta_{A}$ channel and
\[
  \mathrm{IF}_{i}(F_{A})
  = \Bigl(\tfrac{\partial F_{A}}{\partial u}
    + \tfrac{\partial F_{A}}{\partial W}\,DW(u_{0})\Bigr)g_{i}
  = \bigl(-W\rs^{A}+\tfrac12 b^{A}\bigr)\T g_{i},
\]
using $\partial F_{A}/\partial u=-r\T W$ and $\partial(\tfrac12 r\T W r)/\partial
W$ contracted with $DW[g_{i}]=\sum_{l}(\partial W/\partial u_{l})g_{i,l}$.
Subtracting $\mathrm{IF}_{i}(F_{B})$ gives \eqref{eq:cvec}; the central limit
theorem gives the normal law.
\end{proof}

The term $-W(\rs^{A}-\rs^{B})$ is the difference of weighted residuals that a
fixed-weight account would give. The term $\tfrac12(b^{A}-b^{B})$ is the weight
channel, quadratic in each model's own residual and contracted with the weight
derivative. It is the testing analogue of $\Wr$ and is the new ingredient. Every
quantity is already available: $\rs^{A},\rs^{B}$ from the two fits, $\partial
W/\partial u$ from the same derivative the standard error uses, and $\Gamma$ from
the stage-one moments. The relative-fit test is therefore a corollary of the
standard-error machinery, with no projection involved because the envelope
theorem has removed the parameter channel.

\section{The nested case does not commute}
\label{sec:nested}

When $B$ is nested in $A$ the discrepancy difference degenerates under the null
that the restriction holds, $\rs^{A}=\rs^{B}$ in the relevant directions, so
$c=0$ and $\omega^{2}=0$. The first-order term vanishes and the limiting law is
the second-order quadratic form, which is where the projection returns. The
question is whether that quadratic form is Satorra's $U$-difference with the meat
$\Gamma$ simply replaced by the corrected $\Omega$.

\begin{proposition}
\label{prop:nested}
Off the goodness-of-fit null the fitted residual fluctuation is
\[
  \sqrt{N}\,(r_{A}(\widehat\theta_{A})-\rs^{A})
  = -\widetilde M_{A}\sqrt{N}(u-u_{0})+o_{p}(1),
  \qquad
  \widetilde M_{A}=I-\Delta_{A}H_{A}^{-1}\Delta_{A}\T(W-\Wr^{A}),
\]
with $H_{A}=A_{A}+C_{A}$. The nested difference statistic converges to
$\sum_{j}\lambda_{j}\chi^{2}_{1,j}$ with $\lambda_{j}=\operatorname{eig}\!\bigl[
(\widetilde M_{A}\T W\widetilde M_{A}-\widetilde M_{B}\T W\widetilde M_{B})
\Gamma\bigr]$. The generalised annihilator $\widetilde M_{A}$ carries both the
bread curvature $C_{A}$ and the weight channel $\Wr^{A}$ inside the projector, so
it does not equal the fixed-weight annihilator $M_{A}=I-\Delta_{A}A_{A}^{-1}
\Delta_{A}\T W$. Consequently the law is not $\operatorname{eig}[(U_{B}-U_{A})
\Omega]$, and the substitution $\Gamma\mapsto\Omega$ into the $U$-matrix is
incorrect.
\end{proposition}

\begin{proof}
Expand $r_{A}(\widehat\theta_{A})=\rs^{A}+\Delta_{A}(\widehat\theta_{A}-\ths^{A})
-(u-u_{0})+o_{p}(N^{-1/2})$ and substitute the corrected influence
$\sqrt N(\widehat\theta_{A}-\ths^{A})=H_{A}^{-1}\Delta_{A}\T(W-\Wr^{A})\sqrt N
(u-u_{0})+o_{p}(1)$ from \eqref{eq:recap}. Collecting terms gives the stated
$\widetilde M_{A}$. The second-order term of $N(F_{B}-F_{A})$ is the difference
of the quadratic forms $z\T\widetilde M_{A}\T W\widetilde M_{A}z$ and
$z\T\widetilde M_{B}\T W\widetilde M_{B}z$ with $z\sim\mathcal N(0,\Gamma)$, whence
the eigenvalue characterisation. The fixed-weight reduction sets $C_{A}=0$ and
$\Wr^{A}=0$, returning $\widetilde M_{A}=M_{A}$ and $\lambda_{j}=
\operatorname{eig}[(U_{B}-U_{A})\Gamma]$ of \cite{satorra2000}.
\end{proof}

The correction does not factor out of the projection. The weight channel enters
the same place as the bread curvature, inside the annihilator, so the meat and
the projector are entangled. Computing the nested reference law requires the
generalised annihilators $\widetilde M$ rather than a meat swap, and that is a
separate derivation, not a corollary.

\section{Specialization: complete-data DWLS on a linear covariance model}
\label{sec:dwls}

The curvature term $C=\sum_{k}(W\rs)_{k}\nabla^{2}\eta_{k}(\ths)$ is the only
place the shape of the model surface enters. It vanishes whenever $\eta$ is
linear in $\theta$, which removes one of the two deformations and isolates the
weight channel.

\subsection{A model with $C=0$}
Take a linear covariance structure $\Sigma(\theta)=\sum_{k}\theta_{k}G_{k}$ with
fixed symmetric $G_{k}$, so $\eta(\theta)=\vech\Sigma(\theta)=\Delta\theta$ with
$\Delta=[\vech G_{1},\dots,\vech G_{q}]$ constant, $\nabla^{2}\eta_{k}=0$, hence
$C=0$ and $H=A=\Delta\T W\Delta$ constant. A concrete nested pair on $p$
continuous variables: $B$ is compound symmetry (one variance, one covariance),
nested in $A$ (one variance, free covariances), the difference testing equality
of the covariances. Both are linear, so the only surviving deformation is the
weight channel.

\begin{remark}[A test bed, not the headline model]
The parameter is restricted to the open cone $\{\theta:\Sigma(\theta)\succ0\}$;
for compound symmetry this is the wedge $a-b>0$ and $a+(p-1)b>0$ in
$\theta=(a,b)=(\text{variance},\text{covariance})$. This is an inequality
constraint, inactive at an interior pseudo-true value, so it does not enter the
first-order conditions or the asymptotics, and $C=0$ holds from the linearity of
$\eta$ regardless of it. The only design requirement is to keep the
data-generating $\theta_{*}$ interior, with no near-singular best fit. Linear
covariance structures are a classical class (variance components, exchangeable
and Toeplitz patterns, fixed-loading factor models), but pairing one with DWLS is
off the beaten path. Its role here is a curvature-free check: a deformation that
appears here cannot be curvature in disguise. The substantive case keeps a
genuine factor model and carries $C\neq0$ (Section~\ref{sec:cfa}).
\end{remark}

\subsection{DWLS ingredients}
The stage-one moments are $u=\vech(S)$ with casewise influence $g_{i}=\vech(z_{i}
z_{i}\T)-\sigma_{0}$, $z_{i}=y_{i}-\mu$, and $\Gamma=\E[g_{i}g_{i}\T]$ the Browne
fourth-moment matrix. DWLS uses $\widehat W=\diag(\widehat\Gamma)^{-1}$, so with
$\gamma=\diag(\Gamma)$ the population weight is $W=\diag(\gamma)^{-1}$. The weight
influence acts diagonally,
\[
  \dot w_{i,k}=-\gamma_{k}^{-2}\,\dot\gamma_{i,k},\qquad
  \dot\gamma_{i,k}=g_{i,k}^{2}-\gamma_{k}+(\text{centring corrections}),
\]
where $\dot\gamma_{i,k}$ is the casewise influence of the $k$-th diagonal of
$\widehat\Gamma$. The corrected casewise moment-space influence is
\begin{equation}
  v_{i,k}=\frac{g_{i,k}}{\gamma_{k}}
          +\frac{r_{*,k}\,\dot\gamma_{i,k}}{\gamma_{k}^{2}},
  \label{eq:dwls-v}
\end{equation}
which matches the implemented DWLS adapter.

\begin{remark}[The DWLS weight channel is not a matrix on $g_{i}$]
$\widehat W=\diag(\widehat\Gamma)^{-1}$ depends on fourth-order sample moments,
not on the modelled $u=\vech(S)$. Its influence $\dot\gamma_{i,k}$ is quadratic
in $g_{i}$ (through $g_{i,k}^{2}$), so the weight channel in \eqref{eq:dwls-v} is
a genuine casewise quantity, not $\Wr g_{i}$ for a matrix $\Wr$. The matrix
annihilator $\widetilde M$ of Proposition~\ref{prop:nested} therefore holds only
when the weight is a function of $u$; for DWLS the residual influence
$\rho_{i}=\Delta A^{-1}\Delta\T v_{i}-g_{i}$ stays casewise and carries a part
quadratic in $g_{i}$, so its covariance reaches eighth-order moments of the data.
\end{remark}

\subsection{The verification-friendly twin: normal-theory weight}
The normal-theory weight $W=W_{\mathrm{NT}}(S)$ is a function of $u$, so it is the
clean case. With $C=0$ and a data-estimated weight, $\Wr$ is a constant matrix,
$\widetilde M=I-\Delta A^{-1}\Delta\T(W-\Wr)$ is constant, and the nested law is
\[
  \lambda_{j}=\operatorname{eig}\!\bigl[(\widetilde M_{B}\T W\widetilde M_{B}
  -\widetilde M_{A}\T W\widetilde M_{A})\Gamma\bigr],
\]
all constant matrices. A fixed weight would give $\Wr=0$ and, with $C=0$, the
undeformed Satorra law, so the data-estimated weight is exactly what produces the
phenomenon. This isolates the weight channel from the curvature and is the
easiest case to verify by hand and to eigendecompose.

\subsection{The geometry of the payoff}
The deformation is $\widetilde M-M=\Delta A^{-1}\Delta\T\Wr$, linear in $\rs$ and
zero at $\rs=0$. The nested statistic sees $\rs$ only through the rank-$(\mathrm{
df}_{B}-\mathrm{df}_{A})$ map $\widetilde M_{B}-\widetilde M_{A}$ onto the
constraint directions. The exploration is to sweep the direction of $\rs$ at
fixed $\lVert\rs\rVert$: generate from $\Sigma_{0}=\Sigma_{B}(\theta_{*})+
\epsilon D$ and rotate $D$ from $W$-orthogonal to the constraint space toward
aligned. The conjecture, which the simulation tests rather than assumes, is that
orthogonal $\rs$ leaves the law near Satorra (the benign regime where robust
intervals already calibrate) and aligned $\rs$ deforms it most. Because the
deformation is analytic in $\rs$, the predicted miscalibration can overlay the
simulated rejection rate.

\subsection{The deformation in closed form}
\label{sec:closedform}
With $C=0$ the deformation has an exact form, not merely a perturbative one.

\begin{proposition}[Closed-form deformation under $C=0$]
\label{prop:closed}
Let $H=A=\Delta\T W\Delta$ and $\widetilde M=I-\Delta A^{-1}\Delta\T(W-\Wr)$.
Then
\[
  \widetilde M\T W\widetilde M = U + R,\qquad
  U = M\T W M,\qquad
  R = (\Delta\T\Wr)\T A^{-1}(\Delta\T\Wr)\succeq 0,
\]
the cross terms vanishing by the annihilator orthogonality $\Delta\T W M=0$.
Hence the nested difference law is
$\lambda=\operatorname{eig}[(U_{B}-U_{A}+R_{B}-R_{A})\Gamma]$ with mean
$\mathrm{df}_{B}-\mathrm{df}_{A}+\tr[(R_{B}-R_{A})\Gamma]$.
\end{proposition}

\begin{proof}
Write $\widetilde M=M+\delta$ with $\delta=\Delta A^{-1}\Delta\T\Wr\in\operatorname
{col}(\Delta)$. Since $M\T W\Delta=W\Delta-W\Delta A^{-1}(\Delta\T W\Delta)=0$,
both cross terms $M\T W\delta$ and $\delta\T WM$ vanish, and $\delta\T W\delta=
\Wr\T\Delta A^{-1}(\Delta\T W\Delta)A^{-1}\Delta\T\Wr=(\Delta\T\Wr)\T A^{-1}
(\Delta\T\Wr)=R$.
\end{proof}

Two consequences. First, $\Wr$ is linear in $\rs$, so $R$ is quadratic in $\rs$
and the entire deformation is $O(\lVert\rs\rVert^{2})$. The reference law moves at
second order in the base misfit, which is why a benign (small) misspecification
leaves the test near nominal: the first-order effect cancels exactly. Second,
when a model fits, $\rs=0$ gives $R=0$. In a power scenario the larger model fits
($R_{A}=0$), so the mean can only increase, $\mathrm{df}_{B}-\mathrm{df}_{A}+\tr
[R_{B}\Gamma]$; under a true null with a misspecified base both $R$ are present
and the mean shift $\tr[(R_{B}-R_{A})\Gamma]$ is a competition. The numerical
check (\texttt{relative-and-nested-fit-check.py}) confirms the identity to
machine precision and both regimes.

\begin{remark}[Why this is special to $C=0$]
The cancellation uses $A=\Delta\T W\Delta$. With curvature the bread is $H=A+C$,
$M\T W\Delta=W\Delta-W\Delta H^{-1}A\neq0$, and the cross terms survive. The clean
$U+R$ split is exactly what the entanglement of Proposition~\ref{prop:nested}
breaks; the linear test bed buys a closed form that the factor model does not have.
\end{remark}

\subsection{The substantive twin: a factor model with $C\neq0$}
\label{sec:cfa}
Dropping linearity costs only one extra term. For a confirmatory factor model
$\Sigma(\theta)=\Lambda\Psi\Lambda\T+\Theta$ the moments are quadratic in the
loadings, so $\nabla^{2}\eta_{k}$ is sparse and closed form (the $\lambda_{a}
\lambda_{b}$ blocks), and $H=A+C$ with $C=\sum_{k}(W\rs)_{k}\nabla^{2}\eta_{k}
(\ths)$ evaluated at the pseudo-true point. Everything else is unchanged: the
DWLS weight channel \eqref{eq:dwls-v}, the casewise residual influence
$\rho_{i}=\Delta H^{-1}\Delta\T v_{i}-g_{i}$ (now with $H$ in place of $A$), and
the spectrum are assembled exactly as before. A natural nested pair is configural
versus metric invariance across two groups, the difference testing equality of a
block of loadings. The price of realism is that the deformation now superposes
two channels, the curvature $C$ and the weight $\Wr$, so the clean attribution of
the linear test bed is lost. The linear model certifies that the weight channel
alone produces a nonzero effect; the factor model shows the size that effect
reaches in a model people actually fit.

\section{Verdict}
\label{sec:verdict}

The testing extension partitions cleanly. The goodness-of-fit test needs no
correction (Proposition~\ref{prop:delin}); the prevailing scaled statistic is
already right. The relative-fit test for distinguishable models is a corollary
of the standard-error machinery, with $\omega^{2}=c\T\Gamma c$ and the weight
channel \eqref{eq:cvec} computable from quantities the standard error already
forms (Proposition~\ref{prop:vuong}). The nested difference test, in the regime
that does not assume the larger model correct, is its own project, because the
weight channel deforms Satorra's projection rather than its meat
(Proposition~\ref{prop:nested}). For the manuscript this argues to keep the
estimated-weight standard error and the distinguishable relative-fit test
together, and to treat the nested weighted-$\chi^{2}$ law as a sequel
\cite{vuong1989,merkleyou2016}.

\begin{thebibliography}{9}

\bibitem{satorra1989} Satorra, A. (1989). Alternative test criteria in
  covariance structure analysis: a unified approach. \emph{Psychometrika},
  54(1), 131--151.
  \href{https://doi.org/10.1007/BF02294453}{doi:10.1007/BF02294453}.

\bibitem{satorra2000} Satorra, A. (2000). Scaled and adjusted restricted tests
  in multi-sample analysis of moment structures. In R. D. H. Heijmans, D. S. G.
  Pollock, \& A. Satorra (Eds.), \emph{Innovations in Multivariate Statistical
  Analysis} (pp. 233--247). Springer.
  \href{https://doi.org/10.1007/978-1-4615-4603-0_17}{doi:10.1007/978-1-4615-4603-0\_17}.

\bibitem{satorrabentler2001} Satorra, A., \& Bentler, P. M. (2001). A scaled
  difference chi-square test statistic for moment structure analysis.
  \emph{Psychometrika}, 66(4), 507--514.
  \href{https://doi.org/10.1007/BF02296192}{doi:10.1007/BF02296192}.

\bibitem{vuong1989} Vuong, Q. H. (1989). Likelihood ratio tests for model
  selection and non-nested hypotheses. \emph{Econometrica}, 57(2), 307--333.
  \href{https://doi.org/10.2307/1912557}{doi:10.2307/1912557}.

\bibitem{merkleyou2016} Merkle, E. C., You, D., \& Preacher, K. J. (2016).
  Testing nonnested structural equation models. \emph{Psychological Methods},
  21(2), 151--163.
  \href{https://doi.org/10.1037/met0000038}{doi:10.1037/met0000038}.

\bibitem{laisimoes2023} Lai, K., \& Sim\~oes, F. (2023). Reflecting on the
  ``robust'' standard errors for two-stage structural equation modeling.
  \emph{Structural Equation Modeling}, 30(5).

\bibitem{hallinoue2003} Hall, A. R., \& Inoue, A. (2003). The large sample
  behaviour of the generalized method of moments estimator in misspecified
  models. \emph{Journal of Econometrics}, 114(2), 361--394.
  \href{https://doi.org/10.1016/S0304-4076(03)00075-0}{doi:10.1016/S0304-4076(03)00075-0}.

\end{thebibliography}

\end{document}
